% Part: first-order-logic
% Chapter: natural-deduction
% Section: soundness-identity

\documentclass[../../../include/open-logic-section]{subfiles}

\begin{document}

\olfileid[id]{fol}{ntd}{sid}

\olsection{Kesahihan dengan \usetoken{S}{identity}}

\begin{prop}
Deduksi alami dengan aturan untuk $\eq$ bersifat sahih.
\end{prop}

\begin{proof}
Kita memperluas pembuktian induktif teorema kesahihan dengan memeriksa
kasus-kasus untuk dua aturan baru.

Setiap !!{formula} berbentuk $\eq[t][t]$ valid, sebab untuk setiap
!!{structure}~$\Struct M$, berlaku $\Sat{M}{\eq[t][t]}$. (Perhatikan
bahwa kita mengandaikan suku $t$ tertutup, yakni tidak memuat variabel,
sehingga penugasan variabel tidak relevan.)

Untuk varian pertama, andaikan inferensi terakhir dalam !!a{derivation}
adalah \Elim{\eq}. Dengan demikian, !!{derivation} tersebut berbentuk
sebagai berikut:
\begin{prooftree}
  \AxiomC{$\Gamma_1$}
  \RightLabel{$\delta_1$}
  \DeduceC{$\eq[t_1][t_2]$}
  \AxiomC{$\Gamma_2$}
  \RightLabel{$\delta_2$}
  \DeduceC{$!A(t_1)$}
  \RightLabel{\Elim{\eq}}
  \BinaryInfC{$!A(t_2)$}
\end{prooftree}
Kedua premis $\eq[t_1][t_2]$ dan $!A(t_1)$ masing-masing dapat kita
!!{derive} dari asumsi !!{undischarged}~$\Gamma_1$ dan $\Gamma_2$. Kita
hendak menunjukkan bahwa
$!A(t_2)$ merupakan konsekuensi dari $\Gamma_1 \cup \Gamma_2$. Perhatikan
!!a{structure}~$\Struct{M}$ dengan $\Sat{M}{\Gamma_1 \cup \Gamma_2}$.
Menurut hipotesis induksi, $\Sat{M}{!A(t_1)}$ dan
$\Sat{M}{\eq[t_1][t_2]}$. Oleh karena itu,
$\Value{t_1}{M} = \Value{t_2}{M}$. Misalkan $s$ adalah sebarang penugasan
variabel dan $m = \Value{t_1}{M} = \Value{t_2}{M}$. Menurut
\olref[fol][syn][ext]{prop:ext-formulas}, $\Sat{M}{!A(t_1)}[s]$ jika dan
hanya jika $\Sat{M}{!A(x)}[\Subst{s}{m}{x}]$ jika dan hanya jika
$\Sat{M}{!A(t_2)}[s]$. Karena $\Sat{M}{!A(t_1)}$, kita memperoleh
$\Sat{M}{!A(t_2)}$. Varian kedua, dengan arah substitusi sebaliknya,
ditangani dengan cara yang sama.
\end{proof}

\end{document}
